% =====================================================================
% Introduction - material sources (v7.99 / P0-canonical):
%   * framework/proof_status.md sec 3 (claim registry), sec 5 (honest rate)
%   * framework/journal_submission_d1_review.md sec 1.1 (R1), 1.3 (C1-C6)
%   * framework/sorry_contamination_audit.md sec 1 (reversal finding)
%   * framework/ci_slimming_phase1_report.md (batch scale, reproducibility)
%   * framework/p0_reconciliation_report.md (canonical numbers, landed
%     2026-08-21: 196 = 30+44+31+29+30+32; registry 448; residual 253;
%     five classes 50/90/17/24/13 + 2 bundling; Agda 131)
% Number policy: P0-canonical figures appear as plain numbers; the
% remaining \pnum{...} red-bracket values are verified v7.98/v7.99
% working-tree baselines, each annotated with its repository source.
% =====================================================================

\subsection{Setting: a library that scaled faster than its governance}
\label{sec:setting}
TOE-SYLVA is a formalization program whose Lean~4 side is built
strictly on top of a pinned mathlib4 revision (commit
\texttt{8a178386} at the v7.99 baseline; toolchain Lean 4.29.0) and
contains two very different kinds of content. The first is a set of
handcrafted developments targeting physics-adjacent mathematics
(Chern--Simons theory, topological invariants, gauge theory,
$p$-adic analysis, a course chain in algebraic number theory), grown
to roughly \pnum{571611} declared theorems. The second is a
machine-generated \emph{batch corpus} of \pnum{119859} template
modules produced by a generator whose manifest lists \pnum{119831}
entries with SHA-256 sampling (\pnum{500} draws, zero mismatches at
the v7.98 baseline). A third, smaller Agda track holds 7 files with
131 name-level \texttt{postulate} declarations (canonical, P0). The
library, its audits, and this paper were all produced with AI
assistance under human tasking and review; that fact is not a
disclaimer but the paper's subject matter. The defining event of the
program is that the second kind of content --- and, earlier, large
parts of the first --- scaled two to three orders of magnitude faster
than the processes that were supposed to keep the corpus honest.

The result was not one problem but three, each a distinct failure
mode of \emph{measurement}, not of taste:

\begin{enumerate}
\item[\textbf{F1}] \textbf{Axioms as invisible deferred proof.}
An \texttt{axiom} command in Lean is a legal, type-checked, fully
trusted placeholder. In a corpus where AI-assisted drafting produced
plausible-looking statement blocks at scale, the non-batch corpus
accumulated 448 registered \texttt{axiom} declarations (canonical
registry caliber, 116 non-batch files), while top-level summaries
continued to say ``formalized''. ``Formalized'' and ``proved'' are
different predicates; the library had been quietly conflating them.
\item[\textbf{F2}] \textbf{Proof-rate metrics that drift with the
measuring rule.} Over the program's history the reported
``proof rate'' of the handcrafted corpus took the values
\pnum{88.08\%}, \pnum{99.5\%}, \pnum{99.9\%}, \pnum{100\%},
\pnum{99.79\%}, and 3.91\% --- a 25$\times$ spread --- with no change
to the underlying proofs. Every value was ``correct'' under some
counting rule; no rule was stated.
\item[\textbf{F3}] \textbf{String-level scans that mistake
negations for placeholders.} A repository-wide text search for
\texttt{sorry} reported \pnum{30.2\%} of the batch corpus as
``containing sorry'' (\pnum{36199} of \pnum{119859} files). A
full-population code-level audit found the code-level count to be
zero: 36{,}198 of the 36{,}199 hits were the fixed header comment
\emph{asserting the absence} of placeholders, and the single
remaining hit was a theorem whose \emph{name} contains the word.
The scan was measuring comments, not proofs.
\end{enumerate}

None of these failures is exotic. Each is what happens when a metric
is quoted without the rule that produced it --- a \emph{caliber}, in
the terminology this paper develops --- and each would have survived
casual review indefinitely, because the surface numbers were
plausible. They were caught only because the program eventually
built governance whose explicit job was to catch them.

\subsection{The governance program}
\label{sec:program}
The response, executed over versions v7.94--v7.99 of the repository,
is the subject of this paper. It has four stages, which are also the
paper's methodological mainline:

\begin{enumerate}
\item \textbf{Registry audit} (Sec.~\ref{sec:registry}): every
\texttt{axiom} declaration in the non-batch corpus is registered in a
machine-readable JSON registry, classified into a four-class taxonomy
(primitive / definitional / schema / placeholder), with per-entry
provenance. The registry is the single authority for \emph{what the
library owes}: 448 entries at baseline.
\item \textbf{Redemption sweeps} (Sec.~\ref{sec:sweeps}): six
organized passes (sweep1--sweep6) convert registered axioms into
theorems, definitions, or explicitly conditional statements, closing
196 entries in total ($30+44+31+29+30+32$). Every closure is
re-graded \emph{post hoc} by semantic content into a five-class
taxonomy --- vacuous placeholder closure (50), $P\!\to\!P$
conditionalization (90), data-to-definition conversion (17),
definitional reconstruction by \texttt{rfl} (24), substantive proof
(13), plus 2 axiom-bundling restructurings excluded from the
redemption count --- because the honest headline of the program is
\emph{what kind} of work those 196 closures actually were.
\item \textbf{Three-caliber reconciliation}
(Sec.~\ref{sec:reconciliation}): the program's own numbers are then
audited against each other. The registry caliber (448), the live
\texttt{.lean}-declaration caliber (253 residual), and the
sweep-accounting caliber (196 cumulative) close to the identity
$445 - 166 - 26 = 253$ at the declaration level, and
$448 - 253 = 195 = 142 + 42 + 11$ at the registry level. Six
documented caliber conflicts (C1--C6) are resolved and their
anatomy tabulated.
\item \textbf{Honesty governance}
(Secs.~\ref{sec:cschain}--\ref{sec:governance}): a claim-tier system
(THEOREM / THEOREM* / CLAIM / CONJECTURE) with falsifiability
criteria, an honest proof-rate metric defined as a
(number,~caliber) pair, a machine-checked honesty CI, a quarantine
protocol for the one genuinely contaminated module, and a
deleted-claims ledger enforcing the no-silent-deletion rule.
\end{enumerate}

Two lifecycle case studies exercise the whole stack end to end ---
the Chern--Simons chain, where a CLAIM$\to$THEOREM upgrade is
performed along a six-level dependency chain with explicit
mathlib-delegation labels and an explicit statement of what the
upgraded theorem is \emph{not} (Sec.~\ref{sec:cschain}); and the
Dedekind delegation, where a ``theorem'' that is a pure wrapper over
a mathlib instance is registered honestly as such
(Sec.~\ref{sec:dedekind}).

\subsection{Contributions}
\label{sec:contributions}
\begin{itemize}
\item[\textbf{C1}] A four-class axiom taxonomy (primitive /
definitional / schema / placeholder) with a machine-readable registry
of 448 axioms over 116 non-batch files --- the audit artifact that
made F1 visible at all, and the precondition for every later number
in the paper.
\item[\textbf{C2}] A redemption methodology --- six sweeps closing
196 registered axioms --- together with a five-class semantic
re-classification (50 vacuous / 90 $P\!\to\!P$ / 17 data-def / 24
definitional / 13 substantive / 2 bundling excluded) that refuses the
aggregate ``196 redeemed'' reading. The transferable claim is the
\emph{decomposition}, not the count.
\item[\textbf{C3}] A three-caliber reconciliation protocol that
closes the registry (448), live-declaration (253), and sweep (196)
accounts to a machine-reproducible identity, resolving six
documented conflicts (C1--C6) --- including a 25$\times$
proof-rate disagreement and a 2.7$\times$ Agda-postulate
disagreement --- without editing a single proof.
\item[\textbf{C4}] An honesty-governance layer: claim tiers with
falsifiability criteria; an honest proof-rate metric whose reference
implementation counts code-level placeholders only
(\texttt{count\_\allowbreak sorry\_\allowbreak no\_\allowbreak comments}); a repository CI that enforces
the honesty layer's invariants; and the full-population reversal of
the 30.2\% ``sorry contamination'' signal (F3), which we argue is a
\emph{positive} result: the caliber discipline caught a would-be
crisis that string-level tooling had manufactured.
\item[\textbf{C5}] Two full lifecycle case studies --- the
Chern--Simons T3 upgrade with mathlib-delegation labels, gap
accounting, and a concurrent-work rollback that the audit layer
itself detected; and the Dedekind delegation, an honest registration
of zero-new-math ``theorems''. Together they show the tier labels
doing real communicative work.
\end{itemize}

\subsection{Paper roadmap}
\label{sec:roadmap}
Section~\ref{sec:related} positions the work against large-scale
formalization practice (mathlib, Isabelle/AFP, Metamath) and the
2024--2026 literature on LLM-generated proofs and autoformalization
integrity. Sections~\ref{sec:registry}--\ref{sec:reconciliation} are
the methodological core: registry audit, redemption sweeps and the
five-class taxonomy, and the three-caliber reconciliation.
Sections~\ref{sec:cschain} and~\ref{sec:dedekind} are the case
studies. Section~\ref{sec:governance} assembles the honesty layer.
Section~\ref{sec:discussion} states limitations --- including an
explicit disclosure that 71.4\% of closures (the vacuous and
$P\!\to\!P$ classes) transfer no proof content --- anticipated
objections, and future work. Appendix~A gives the reproduction
protocol; Appendix~B is a glossary of calibers.

\subsection{Transparency statement}
\label{sec:transparency}
This library was produced with AI assistance, and so was this paper.
The repository discipline it describes --- registration first,
boundaries explicit, no silent deletions --- is applied to the
manuscript itself. All headline figures (196, 448, 253, the
five-class decomposition, Agda 131) are canonical per the P0
reconciliation report
(\texttt{framework/p0\_reconciliation\_report.md}, 2026-08-21~\cite{sylvap0recon2026}), and
the remaining values in red angle brackets are declared
v7.98/v7.99 working-tree baselines, each annotated with its
repository source in the \LaTeX{} source comments. Every claim about
the library cites the internal audit report it rests on. The
supplementary build artifacts (compile log, PDF) accompany the
manuscript under \texttt{framework/paper/}.
